% Options for packages loaded elsewhere
\PassOptionsToPackage{unicode}{hyperref}
\PassOptionsToPackage{hyphens}{url}
%
\documentclass[
]{article}
\usepackage{amsmath,amssymb}
\usepackage{iftex}
\ifPDFTeX
  \usepackage[T1]{fontenc}
  \usepackage[utf8]{inputenc}
  \usepackage{textcomp} % provide euro and other symbols
\else % if luatex or xetex
  \usepackage{unicode-math} % this also loads fontspec
  \defaultfontfeatures{Scale=MatchLowercase}
  \defaultfontfeatures[\rmfamily]{Ligatures=TeX,Scale=1}
\fi
\usepackage{lmodern}
\ifPDFTeX\else
  % xetex/luatex font selection
\fi
% Use upquote if available, for straight quotes in verbatim environments
\IfFileExists{upquote.sty}{\usepackage{upquote}}{}
\IfFileExists{microtype.sty}{% use microtype if available
  \usepackage[]{microtype}
  \UseMicrotypeSet[protrusion]{basicmath} % disable protrusion for tt fonts
}{}
\makeatletter
\@ifundefined{KOMAClassName}{% if non-KOMA class
  \IfFileExists{parskip.sty}{%
    \usepackage{parskip}
  }{% else
    \setlength{\parindent}{0pt}
    \setlength{\parskip}{6pt plus 2pt minus 1pt}}
}{% if KOMA class
  \KOMAoptions{parskip=half}}
\makeatother
\usepackage{xcolor}
\usepackage[margin=1in]{geometry}
\usepackage{graphicx}
\makeatletter
\newsavebox\pandoc@box
\newcommand*\pandocbounded[1]{% scales image to fit in text height/width
  \sbox\pandoc@box{#1}%
  \Gscale@div\@tempa{\textheight}{\dimexpr\ht\pandoc@box+\dp\pandoc@box\relax}%
  \Gscale@div\@tempb{\linewidth}{\wd\pandoc@box}%
  \ifdim\@tempb\p@<\@tempa\p@\let\@tempa\@tempb\fi% select the smaller of both
  \ifdim\@tempa\p@<\p@\scalebox{\@tempa}{\usebox\pandoc@box}%
  \else\usebox{\pandoc@box}%
  \fi%
}
% Set default figure placement to htbp
\def\fps@figure{htbp}
\makeatother
\setlength{\emergencystretch}{3em} % prevent overfull lines
\providecommand{\tightlist}{%
  \setlength{\itemsep}{0pt}\setlength{\parskip}{0pt}}
\setcounter{secnumdepth}{-\maxdimen} % remove section numbering
\usepackage{booktabs}
\usepackage{caption}
\usepackage{longtable}
\usepackage{colortbl}
\usepackage{array}
\usepackage{anyfontsize}
\usepackage{multirow}
\usepackage{bookmark}
\IfFileExists{xurl.sty}{\usepackage{xurl}}{} % add URL line breaks if available
\urlstyle{same}
\hypersetup{
  pdftitle={Analyse du risque de récidive du syndrome oculo-respiratoire après la vaccination contre la grippe ».},
  pdfauthor={Saa Moussa TENGUIANO, Michelle Tchagang NGOUKEU, Biostatistique, Université Laval},
  hidelinks,
  pdfcreator={LaTeX via pandoc}}

\title{Analyse du risque de récidive du syndrome oculo-respiratoire
après la vaccination contre la grippe ».}
\author{Saa Moussa TENGUIANO, Michelle Tchagang NGOUKEU, Biostatistique,
Université Laval}
\date{}

\begin{document}
\maketitle

{
\setcounter{tocdepth}{2}
\tableofcontents
}
\textbf{Resumé:} Ce rapport présente une analyse approfondie du risque
de récidive du syndrome ocul-orespiratoire après la vaccination
antigrippale. L'objectif est de comparer le risque de récidive entre les
différents vaccins et le placebo, et de déterminer si ce risque varie
selon le type de vaccin. Pour ce faire, une étude descriptive ainsi
qu'une étude analytique utilisant un modèle GEE ont été réalisées afin
d'estimer le risque de récidive du syndrome oculo-respiratoire après
vaccination.

\section{\texorpdfstring{\emph{Introduction.}}{Introduction.}}\label{introduction.}

\subsection{Contexte.}\label{contexte.}

Le syndrome oculo respiratoire (SOR) est un effet secondaire
généralement bénin de la vaccination antigrippale, se manifestant par
une rougeur des yeux, un mal de gorge, de la toux, des difficultés
respiratoires ou un œdème facial. Ces symptômes apparaissent le plus
souvent entre 2 et 12 heures après la vaccination, mais certains cas
surviennent dès quelques minutes ou jusqu'à 48 à 72 heures après
l'injection (INSPQ, 2001).

Le SOR a été identifié pour la première fois au Québec lors de la
campagne de vaccination contre l'influenza 2000-2001, ce qui a soulevé
des inquiétudes quant à la sécurité des vaccins. La majorité des cas
étaient associés au vaccin Fluviral S/F, tandis qu'une minorité
concernait d'autres vaccins comme Vaxigrip et Fluzone. Pour réduire le
risque, Fluviral S/F a été modifié lors de la campagne 2001-2002, ce qui
a entraîné une diminution significative du nombre de cas et une
atténuation des symptômes, bien que le SOR n'ait pas complètement
disparu (INSPQ, 2001).

La question de la récidive chez les personnes ayant déjà présenté un SOR
a été étudiée par Skowronski et al.~(2003) dans le cadre d'un essai
clinique randomisé et contrôlé par placebo, qui a été interrompu lorsque
le risque attribuable au vaccin est devenu trop élevé. En parallèle, une
enquête téléphonique menée par Skowronski et al.~(2002) a montré que le
risque de récidive après revaccination était estimé à 5 \% (IC à 95\,\%
: 2,2\,\%-10,5\,\%).

Dans ce contexte, l'objectif de l'étude analysée était d'évaluer le
risque de récidive du SOR chez les personnes affectées en 2000 2001,
pour chacun des vaccins disponibles au Canada lors de la saison grippale
2002-2003, afin d'éclairer la pratique de la revaccination et les
recommandations de santé publique.

Ainsi, cette étude illustre de manière concrète plusieurs concepts clés
abordés dans le cours sur les essais cliniques et les interventions,
notamment le choix du devis expérimental (essai clinique randomisé en
double aveugle avec placebo), la définition de la population cible, la
mise en place de l'intervention, la distinction entre issues principales
et secondaires, ainsi que l'analyse et l'interprétation des résultats.

\subsection{Conception de l'étude}\label{conception-de-luxe9tude}

L'étude était un essai clinique contrôlé à plan croisé, mené dans trois
sites au Canada (Vancouver, Québec, Montérégie). Trois groupes de
patients ≥ 18 ans ayant déjà présenté un syndrome oculo-respiratoire
(SOR) ont été constitués selon leur historique vaccinal antérieur. Les
patients déjà vaccinés, indisponibles pour le suivi ou présentant des
conditions pouvant masquer des symptômes du SOR étaient exclus. Chaque
participant recevait deux injections à 7 jours d'intervalle : une dose
de vaccin (Fluviral S/F ou Vaxigrip) et un placebo, selon une séquence
randomisée en blocs de 4. La randomisation était effectuée séparément
par site et par groupe. L'aveugle était maintenu grâce à des seringues
opaques préparées par une infirmière indépendante. Le placebo était une
solution saline isotonique. Les vaccins utilisés provenaient d'un seul
lot chacun et contenaient des souches grippales standardisées, avec
différents agents de fragmentation selon le produit. La collecte des
données comprenait un carnet de suivi quotidien des symptômes (jours 1 à
7), des appels téléphoniques au jour 1 (pour les symptômes initiaux) et
au jour 3 si symptômes, ainsi qu'une visite de suivi environ une semaine
plus tard lors de la seconde injection. Le SOR était défini selon la
définition nationale de 2001 : apparition dans les 24 h post-vaccination
d'au moins un symptôme typique (rougeur oculaire bilatérale, toux,
dyspnée, œdème facial, etc.). Des définitions plus strictes étaient
explorées en analyses secondaires. La sévérité était classée en trois
niveaux selon l'impact sur les activités quotidiennes.

\subsection{Objectifs et plan
d'analyse.}\label{objectifs-et-plan-danalyse.}

L'objectif principal de ce rapport est d'évaluer le risque de récidive
du syndrome oculo-respiratoire après une vaccination antigrippale. Plus
précisément, il s'agit : de comparer le risque de récidive du SOR entre
les vaccins et le placebo, et de comparer les risques de récidive entre
les différents vaccins. Pour atteindre ces objectifs, l'analyse sera
structurée de la manière suivante : une analyse descriptive du risque de
récidive, et un modèle de régression pour évaluer le risque de récidive
en tenant compte des facteurs pertinents.

\section{\texorpdfstring{\emph{Analyse du risque de récidive du syndrome
oculo respiratoire après la vaccination contre la
grippe.}}{Analyse du risque de récidive du syndrome oculo respiratoire après la vaccination contre la grippe.}}\label{analyse-du-risque-de-ruxe9cidive-du-syndrome-oculo-respiratoire-apruxe8s-la-vaccination-contre-la-grippe.}

Afin d'analyser le risque de récidive du syndrome oculo-respiratoire,
nous commencerons par décrire la base de données issue de l'étude et
créer les variables supplémentaires nécessaires à l'analyse. Cette étape
permettra également de vérifier les données manquantes et d'assurer la
cohérence avant de réaliser les analyses descriptives et les modèles
statistiques.

\subsection{Description de la base de
données}\label{description-de-la-base-de-donnuxe9es}

\begin{verbatim}
  [1] 292   8
\end{verbatim}

\begin{table}[t]
\caption*{
{\fontsize{20}{25}\selectfont  Aper\c{c}u des 4 premi\`eres lignes de la base\fontsize{12}{15}\selectfont }
} 
\fontsize{12.0pt}{14.0pt}\selectfont
\begin{tabular*}{\linewidth}{@{\extracolsep{\fill}}rrlrllrr}
\toprule
AGE & SEXE & CATEGORIE & id & region & traitement & dose & SOR:>=1 sympt\^ome:24h \\ 
\midrule\addlinespace[2.5pt]
64 & 1 & B & 100 & M & Placebo-Fluviral & 1 & 2 \\ 
64 & 1 & B & 100 & M & Placebo-Fluviral & 2 & 2 \\ 
46 & 1 & B & 101 & M & Placebo-Vaxigrip & 1 & 2 \\ 
\bottomrule
\end{tabular*}
\end{table}

La base de données comprend \textbf{192 observations} pour 146
participants réparties sur \textbf{8 variables principales}:

\begin{itemize}
\tightlist
\item
  \textbf{age} : âge du participant au moment de l'inclusion.\\
\item
  \textbf{catégorie} : variable factorielle à trois modalités (A, B et
  C) définissant le groupe d'appartenance des participants.

  \begin{itemize}
  \tightlist
  \item
    \textbf{Groupe A} : participants ayant présenté un SOR en 2000-2001
    et \textbf{non revaccinés} contre la grippe.\\
  \item
    \textbf{Groupe B} : participants ayant présenté un SOR en 2000-2001
    et \textbf{revaccinés} en 2001-2002.\\
  \item
    \textbf{Groupe C} : participants présentant une \textbf{première
    occurrence} de SOR en 2001-2002.\\
  \end{itemize}
\item
  \textbf{SOR} : indicateur de la présence d'au moins un symptôme de SOR
  survenant dans les \textbf{24 heures} suivant la vaccination.\\
\item
  \textbf{dose} : ordre des interventions, correspondant à la
  \textbf{première ou la deuxième injection}.\\
\item
  \textbf{identifiant} : identifiant unique attribué à chaque
  participant, chacun ayant \textbf{deux observations} correspondant aux
  deux doses reçues.\\
\item
  \textbf{région} : région de recrutement des participants (Vancouver,
  Québec ou Montérégie).\\
\item
  \textbf{traitement} : intervention administrée à chaque visite,
  précisant \textbf{placebo ou vaccin}, ainsi que le type de vaccin
  utilisé.
\end{itemize}

Pour la suite de l'analyse, deux variables supplémentaires seront créées
:

\textbf{Traitement reçu} : indiquant le traitement effectivement reçu
par le participant en tenant compte de la dose. \textbf{SOR
dichotomique} : variable binaire valant 1 si le participant présente au
moins un symptôme de SOR et 0 sinon.

\subsection{Création des variables}\label{cruxe9ation-des-variables}

\begin{table}[t]
\fontsize{12.0pt}{14.0pt}\selectfont
\begin{tabular*}{\linewidth}{@{\extracolsep{\fill}}rrlrllrrlr}
\toprule
AGE & SEXE & CATEGORIE & id & region & traitement & dose & SOR:>=1 sympt\^ome:24h & traitement\_recu & SOR\_bin \\ 
\midrule\addlinespace[2.5pt]
64 & 1 & B & 100 & M & Placebo-Fluviral & 1 & 2 & Placebo & 0 \\ 
64 & 1 & B & 100 & M & Placebo-Fluviral & 2 & 2 & Fluviral & 0 \\ 
\bottomrule
\end{tabular*}
\end{table}

La nouvelle base de données comporte désormais 10 variables. Pour
l'analyse descriptive du risque de récidive, nous retiendrons les
variables suivantes : catégorie, région, traitement, traitement reçu et
la variable SOR\_bin indiquant la présence ou l'absence de symptômes.

\subsection{Caractéristiques des 146 participants par
groupe}\label{caractuxe9ristiques-des-146-participants-par-groupe}

\begin{table}[t]
\caption*{
{\fontsize{20}{25}\selectfont  R\\'epartition des participants par r\\'egion et groupe\fontsize{12}{15}\selectfont }
} 
\fontsize{12.0pt}{14.0pt}\selectfont
\begin{tabular*}{\linewidth}{@{\extracolsep{\fill}}lrrr}
\toprule
R\\'egion & A & B & C \\ 
\midrule\addlinespace[2.5pt]
M & 33 & 38 & 20 \\ 
Q & 17 & 34 & 68 \\ 
V & 50 & 28 & 12 \\ 
\bottomrule
\end{tabular*}
\end{table}

La répartition géographique des participants varie selon le groupe ORS :
le \textbf{groupe A} est majoritairement recruté à Vancouver et
Montérégie, avec une faible représentation à Québec ; le \textbf{groupe
B} est réparti de manière relativement équilibrée entre Montérégie et
Québec, avec moins de participants à Vancouver ; enfin, le
\textbf{groupe C} est principalement recruté à Québec, avec une présence
beaucoup plus faible à Montérégie et Vancouver.

\subsection{\texorpdfstring{\emph{Comparaison des risques de récidive du
syndrome oculo-respiratoire entre les vaccins et le
placébo}}{Comparaison des risques de récidive du syndrome oculo-respiratoire entre les vaccins et le placébo}}\label{comparaison-des-risques-de-ruxe9cidive-du-syndrome-oculo-respiratoire-entre-les-vaccins-et-le-placuxe9bo}

\subsubsection{Comparaison des taux d'incidence du syndrome
oculo-respiratoire entre le vaccin fluviral et le
placébo}\label{comparaison-des-taux-dincidence-du-syndrome-oculo-respiratoire-entre-le-vaccin-fluviral-et-le-placuxe9bo}

\paragraph{Calcul des proportions
d'incidice}\label{calcul-des-proportions-dincidice}

Le tableau ci-dessous présente le nombre de cas de syndrome
oculo‑respiratoire chez les patients ayant reçu le placebo, comparé à
ceux ayant reçu le Fluviral, ainsi que les proportion d'incidence
correspondants.

\begin{table}[t]
\caption*{
{\fontsize{20}{25}\selectfont  R\\'esum\\'e des cas : Placebo contre Fluviral\fontsize{12}{15}\selectfont }
} 
\fontsize{12.0pt}{14.0pt}\selectfont
\begin{tabular*}{\linewidth}{@{\extracolsep{\fill}}lrrr}
\toprule
Vaccin re\c{c}u & Nombre total & Nombre de cas & Incidence (\%) \\ 
\midrule\addlinespace[2.5pt]
Fluviral & 73 & 31 & 42.5 \\ 
Placebo & 73 & 6 & 8.2 \\ 
\bottomrule
\end{tabular*}
\end{table}

Le tableau montre que la probabilité de récidive du syndrome
oculo‑respiratoire est de 42,5 \% chez les patients ayant reçu le vaccin
Fluviral, contre 8,2 \% chez ceux ayant reçu le placebo.

\paragraph{Calcul de la différence d'incidence et l'intervalle de
confiance}\label{calcul-de-la-diffuxe9rence-dincidence-et-lintervalle-de-confiance}

Le tableau ci‑dessous présente la différence des taux d'incidence entre
le vaccin Fluviral et le placebo.

\begin{table}[t]
\caption*{
{\fontsize{20}{25}\selectfont  Diff\\'erence d'incidence : Fluviral contre Placebo\fontsize{12}{15}\selectfont }
} 
\fontsize{12.0pt}{14.0pt}\selectfont
\begin{tabular*}{\linewidth}{@{\extracolsep{\fill}}lrl}
\toprule
Vaccins & Diff\\'erence d'incidence & Intervalle de confiance \\ 
\midrule\addlinespace[2.5pt]
Fluviral contre Placebo & 34\% & (IC \`a 95\%: 21\%-47\%) \\ 
\bottomrule
\end{tabular*}
\end{table}

La vaccination avec Fluviral augmente la probabilité de récidive du
syndrome oculo‑respiratoire de 34 \% par rapport au placebo (IC 95 \% :
21\%‑47 \%). Comme l'intervalle de confiance n'inclut pas la valeur
nulle, cette différence est statistiquement significative. Cependant,
l'intervalle reste relativement large, ce qui reflète une certaine
incertitude autour de l'estimation.

\subsubsection{Comparaison des taux d'incidence du syndrome
oculo-respiratoire entre le vaccin Vaxigrip et le
placébo}\label{comparaison-des-taux-dincidence-du-syndrome-oculo-respiratoire-entre-le-vaccin-vaxigrip-et-le-placuxe9bo}

\paragraph{Calcul des proportions
d'incidice}\label{calcul-des-proportions-dincidice-1}

Le tableau ci-dessous présente le nombre de cas de syndrome
oculo‑respiratoire chez les patients ayant reçu le placebo, comparé à
ceux ayant reçu le Vaxigrip, ainsi que les proportion d'incidence
correspondants.

\begin{table}[t]
\caption*{
{\fontsize{20}{25}\selectfont  R\\'esum\\'e des cas : Placebo contre Vaxigrip\fontsize{12}{15}\selectfont }
} 
\fontsize{12.0pt}{14.0pt}\selectfont
\begin{tabular*}{\linewidth}{@{\extracolsep{\fill}}lrrr}
\toprule
Vaccin re\c{c}u & Nombre total & Nombre de cas & Incidence (\%) \\ 
\midrule\addlinespace[2.5pt]
Placebo & 73 & 10 & 13.7 \\ 
Vaxigrip & 73 & 21 & 28.8 \\ 
\bottomrule
\end{tabular*}
\end{table}

Le tableau indique que la probabilité de récidive du syndrome
oculo‑respiratoire est de 28,8 \% chez les patients ayant reçu le vaccin
Vaxigrip, contre 13,7 \% chez ceux ayant reçu le placebo.

\paragraph{Calcul de la différence d'incidence et l'intervalle de
confiance}\label{calcul-de-la-diffuxe9rence-dincidence-et-lintervalle-de-confiance-1}

Le tableau suivant présente la différence des taux d'incidence entre le
vaccin Vaxigrip et le placebo.

\begin{table}[t]
\caption*{
{\fontsize{20}{25}\selectfont  Diff\\'erence d'incidence\fontsize{12}{15}\selectfont }
} 
\fontsize{12.0pt}{14.0pt}\selectfont
\begin{tabular*}{\linewidth}{@{\extracolsep{\fill}}lrl}
\toprule
Vaccins & Diff\\'erence d'incidence & Intervalle de confiance \\ 
\midrule\addlinespace[2.5pt]
Vaxigrip vs Placebo & 15\% & (IC \`a 95\%: 2\%-28\%) \\ 
\bottomrule
\end{tabular*}
\end{table}

La vaccination avec Vaxigrip augmente la probabilité de récidive du
syndrome oculo‑respiratoire de 15\,\% par rapport au placebo (IC 95\,\%
: 2\,\%‑28\,\%). Comme l'intervalle de confiance n'inclut pas zéro,
cette différence est statistiquement significative. Néanmoins,
l'intervalle reste relativement large, traduisant une certaine
incertitude autour de l'estimation.

\textbf{Conclusion}: L'analyse descriptive montre que les deux vaccins
étudiés, Vaxigrip et Fluviral, augmentent significativement la
probabilité de récidive du syndrome oculo‑respiratoire par rapport au
placebo, avec des différences d'incidence de 15\,\% et 34\,\%,
respectivement. Dans les deux cas, les intervalles de confiance ne
comprennent pas zéro, indiquant une significativité statistique.
Toutefois, la largeur relative de ces intervalles traduit une certaine
incertitude autour des estimations, suggérant que les effets réels
pourraient varier dans la population.

\subsubsection{Comparaison des taux d'incidence du syndrome
oculo-respiratoire entre les vaccins fluviral et
Vaxigrip}\label{comparaison-des-taux-dincidence-du-syndrome-oculo-respiratoire-entre-les-vaccins-fluviral-et-vaxigrip}

\begin{table}[t]
\caption*{
{\fontsize{20}{25}\selectfont  Diff\\'erence d\textquoterightincidence : Fluviral vs Vaxigrip\fontsize{12}{15}\selectfont }
} 
\fontsize{12.0pt}{14.0pt}\selectfont
\begin{tabular*}{\linewidth}{@{\extracolsep{\fill}}lrl}
\toprule
Vaccins & Diff\\'erence d'incidence & Intervalle de confiance \\ 
\midrule\addlinespace[2.5pt]
Fluviral contre Vaxigrip & 13.7\% & 13.7(IC \`a 95\%: -1.7\%-29.1\%) \\ 
\bottomrule
\end{tabular*}
\end{table}

\textbf{Conclusion}: La différence d'incidence du syndrome
oculo‑respiratoire entre Fluviral et Vaxigrip est de 13,7\,\% (IC 95\,\%
: -1,7\,\% à 29,1\,\%). Comme l'intervalle de confiance inclut zéro,
cette différence n'est pas statistiquement significative, indiquant
qu'aucun vaccin n'a montré un avantage clair sur l'autre dans cette
étude.

\section{Modélisation statistique du risque de récidive du syndrome
oculo-respiratoire}\label{moduxe9lisation-statistique-du-risque-de-ruxe9cidive-du-syndrome-oculo-respiratoire}

Afin d'analyser plus profondement la différence de risque entre les
vaccins Fluviral, Vaxigrip et le Placebo, nous allons utiliser un modèle
statisitque adaptée à la structure de nos données.

\section{Reférences}\label{refuxe9rences}

1-Skowronski, D. M., De Serres, G., Scheifele, D., Russell, M. L.,
Warrington, R., Davies, H. D., Dionne, M., Duval, B., Kellner, J., \&
MacDonald, J. (2003). Randomized, double-blind, placebo-controlled trial
to assess recurrence risk of oculorespiratory syndrome following
influenza vaccination in previously affected persons. Clinical
Infectious Diseases, 37(8), 1059--1066.
\url{https://doi.org/10.1086/378274}

2-Skowronski, D. M., Strauss, B., Kendall, P., Duval, B., \& De Serres,
G. (2002). Low risk of recurrence of oculorespiratory syndrome following
influenza revaccination. CMAJ, 167(8), 853--858.

3-Groupe scientifique en immunisation, Institut national de santé
publique du Québec. (2001, novembre). Syndrome oculo-respiratoire
observé après la vaccination contre l'influenza. INSPQ.
\url{http://www.inspq.qc.ca}

4-De Serres, G., Skowronski, D. M., Guay, M., Rochette, L., Jacobsen,
K., Fuller, T., \& Duval, B. (2003). Recurrence risk of oculorespiratory
syndrome after influenza vaccination: Randomized controlled trial of
previously affected persons. Clinical Infectious Diseases, 37(8),
1059--1066. \url{https://doi.org/10.1086/378274}

\end{document}
